\section{Part IV: Scale Flow and Content}

\subsection{RG}

The renormalisation group arises as a theorem about how couplings must shift to preserve predictions when the fiducial log-scale changes, as required by Axiom A5. This section formalises three main results: that scale flow forms a one-parameter group, that self-similar observables obey power-law evolution (the Callan--Symanzik equation), and that the effective field theory tower of thresholds is forced by the scale structure alone. Each excitation carries a log-threshold, and the content available at any log-energy is exactly the set of excitations whose thresholds lie at or below it; the structure of the renormalisation group is therefore a consequence of scale invariance, not an input.

\begin{definition}
A \textbf{scale flow} on a space $S$ of dimensionless couplings is the structure
\begin{equation}
\mathsf{ScaleFlow}(S) := \{
  \text{flow} : \mathbb{R} \to S \to S
  \mid
  \text{flow}(0, g) = g \text{ and } \text{flow}(t+u, g) = \text{flow}(t, \text{flow}(u, g))
\}
\end{equation}
It represents the action of a shift of the fiducial log-scale on the coupling space. \lean{RG.ScaleFlow}
\end{definition}

\begin{definition}
A coupling $g \in S$ is a \textbf{fixed point} of a scale flow $\Phi$ if $\Phi.\text{flow}(t, g) = g$ for all $t \in \mathbb{R}$. \lean{RG.ScaleFlow.IsFixedPoint}
\end{definition}

\begin{theorem}
At a fixed point of the flow, every observable is constant under the flow. More precisely, if $\Phi.\text{IsFixedPoint}(g)$ holds and $O : S \to \beta$ is any observable, then for all $t \in \mathbb{R}$,
\begin{equation}
O(\Phi.\text{flow}(t, g)) = O(g) .
\end{equation}
\lean{RG.ScaleFlow.observable\_const\_at\_fixedPoint}
\end{theorem}

\begin{proof}
Immediate from the definition of fixed point.
\end{proof}

\begin{theorem}
If $g$ is a fixed point of the scale flow $\Phi$, then $\Phi.\text{flow}(t, g)$ is also a fixed point for any $t \in \mathbb{R}$.
\lean{RG.ScaleFlow.fixedPoint\_invariant}
\end{theorem}

\begin{proof}
For any $u \in \mathbb{R}$, we have $\Phi.\text{flow}(u, \Phi.\text{flow}(t, g)) = \Phi.\text{flow}(t, \Phi.\text{flow}(u, g)) = \Phi.\text{flow}(t, g)$ by the group law and the fixed-point property of $g$.
\end{proof}

\begin{theorem}[Callan--Symanzik equation and self-similarity]
Suppose $O : \mathbb{R} \to \mathbb{R}$ satisfies the Callan--Symanzik equation
\begin{equation}
\frac{dO}{dt}(t) = \gamma \cdot O(t)
\end{equation}
for all $t \in \mathbb{R}$, where $\gamma$ is a constant (the anomalous dimension). Then
\begin{equation}
O(t) = O(0) \cdot e^{\gamma t} .
\end{equation}
In particular, $\gamma = 0$ (a fixed point) yields exact scale invariance: $O(t) = O(0)$. \lean{RG.callan\_symanzik}
\end{theorem}

\begin{proof}
Define $u(r) := O(r) \cdot e^{-\gamma r}$. The derivative of $u$ is
\begin{equation}
\dot{u}(r) = \dot{O}(r) \cdot e^{-\gamma r} + O(r) \cdot (-\gamma e^{-\gamma r})
  = e^{-\gamma r} (\gamma O(r) - \gamma O(r)) = 0 .
\end{equation}
Thus $u$ is constant, so $u(t) = u(0)$, giving $O(t) e^{-\gamma t} = O(0)$, hence $O(t) = O(0) e^{\gamma t}$.
\end{proof}

\begin{theorem}
At a fixed point of the flow ($\gamma = 0$), the observable is exactly constant: $O(t) = O(0)$ for all $t$.
\lean{RG.scale\_invariant\_of\_fixedPoint}
\end{theorem}

\begin{proof}
Follows from the Callan--Symanzik theorem with $\gamma = 0$.
\end{proof}

\begin{definition}
Given a function $\mu : \iota \to \mathbb{R}$ assigning a log-threshold to each element $i \in \iota$ (where $\iota$ is a finite type), and a log-energy $t \in \mathbb{R}$, the \textbf{active content} at that log-energy is
\begin{equation}
\text{active}(\mu, t) := \{i \in \iota : \mu(i) \leq t\} .
\end{equation}
\lean{RG.active}
\end{definition}

\begin{theorem}
An element $i$ is in $\text{active}(\mu, t)$ if and only if $\mu(i) \leq t$.
\lean{RG.mem\_active}
\end{theorem}

\begin{proof}
Immediate from the definition.
\end{proof}

\begin{theorem}
The active content is monotone in the log-energy: if $t \leq t'$, then $\text{active}(\mu, t) \subseteq \text{active}(\mu, t')$. That is, new physics appears as the scale rises, and never disappears.
\lean{RG.active\_mono}
\end{theorem}

\begin{proof}
If $i \in \text{active}(\mu, t)$, then $\mu(i) \leq t \leq t'$, so $i \in \text{active}(\mu, t')$.
\end{proof}

\begin{theorem}
Between thresholds the effective theory is exact. If $t \leq t'$ and no threshold lies in the interval $(t, t']$, that is, $\neg \exists i : t < \mu(i) \leq t'$, then
\begin{equation}
\text{active}(\mu, t) = \text{active}(\mu, t') .
\end{equation}
\lean{RG.active\_eq\_of\_no\_threshold}
\end{theorem}

\begin{proof}
Since the active set is monotone, it suffices to show $\text{active}(\mu, t') \subseteq \text{active}(\mu, t)$. Suppose $i \in \text{active}(\mu, t')$, so $\mu(i) \leq t'$. If $\mu(i) > t$, then $t < \mu(i) \leq t'$, contradicting the hypothesis. Thus $\mu(i) \leq t$, so $i \in \text{active}(\mu, t)$.
\end{proof}

\begin{theorem}
The tower of effective theories is finite. The range $\{\text{active}(\mu, t) : t \in \mathbb{R}\}$ is a finite set of subsets of $\iota$, so there are only finitely many distinct effective theories as $t$ varies. Changing the effective framework at each threshold is therefore not a defect of the method; it is the exact bookkeeping of a scale-stratified world.
\lean{RG.active\_finite\_range}
\end{theorem}

\begin{proof}
The range is a subset of the power set $\mathcal{P}(\iota)$, which is finite since $\iota$ is finite.
\end{proof}

\begin{theorem}
Below all thresholds the effective theory is empty; above all thresholds it is complete.
\begin{enumerate}
\item If $\mu(i) \leq t$ for all $i \in \iota$, then $\text{active}(\mu, t) = \iota$.
\item If $t < \mu(i)$ for all $i \in \iota$, then $\text{active}(\mu, t) = \emptyset$.
\end{enumerate}
\lean{RG.active\_eq\_univ\_of\_ge} and \lean{RG.active\_eq\_empty\_of\_lt}
\end{theorem}

\begin{proof}
(1) Every element $i$ satisfies $\mu(i) \leq t$, so belongs to $\text{active}(\mu, t)$.
(2) No element $i$ satisfies $\mu(i) \leq t$ since $t < \mu(i)$, so $\text{active}(\mu, t) = \emptyset$.
\end{proof}

\begin{definition}
The \textbf{beta coefficient} as a function of the active content is defined by
\begin{equation}
\beta_\Phi(\mu, t) := \Phi(\text{active}(\mu, t)) ,
\end{equation}
where $\Phi : \mathcal{P}(\iota) \to \mathbb{R}$ maps each finite subset of $\iota$ to a real number. This expresses the idea that the compensation rate required by A5 depends only on which excitations are present.
\lean{RG.betaOf}
\end{definition}

\begin{theorem}
The beta coefficient is constant between thresholds. If $t \leq t'$ and no threshold lies in $(t, t']$, then
\begin{equation}
\beta_\Phi(\mu, t) = \beta_\Phi(\mu, t') .
\end{equation}
Where the content does not change, the compensation required by A5 does not change either.
\lean{RG.beta\_const\_between\_thresholds}
\end{theorem}

\begin{proof}
Since $\text{active}(\mu, t) = \text{active}(\mu, t')$ by the preceding theorem, we have $\beta_\Phi(\mu, t) = \Phi(\text{active}(\mu, t)) = \Phi(\text{active}(\mu, t')) = \beta_\Phi(\mu, t')$.
\end{proof}

\begin{theorem}
The beta function takes finitely many values: the set $\{\beta_\Phi(\mu, t) : t \in \mathbb{R}\}$ is finite. It is a step function of the log-scale whose jumps sit exactly at the masses. The framework fixes this structure; it does not fix the numbers, which remain an input.
\lean{RG.beta\_finite\_range}
\end{theorem}

\begin{proof}
The range of $\beta_\Phi(\mu, \cdot)$ is contained in $\Phi(\text{active}(\mu, \mathbb{R})) = \Phi(\text{range}(\text{active}(\mu)))$, which is finite since $\text{active}(\mu)$ has finite range.
\end{proof}

\begin{theorem}
Above every threshold the coefficient settles to its complete-content value. If $\mu(i) \leq t$ for all $i$, then
\begin{equation}
\beta_\Phi(\mu, t) = \Phi(\iota) .
\end{equation}
The deep-ultraviolet running is governed by the full spectrum.
\lean{RG.beta\_eq\_of\_all\_active}
\end{theorem}

\begin{proof}
If all thresholds are below $t$, then $\text{active}(\mu, t) = \iota$, so $\beta_\Phi(\mu, t) = \Phi(\iota)$.
\end{proof}

\subsection{Running}

In a scale-covariant world, no log-scale is preferred, so the defects capable of being resolved by a probe cannot bunch anywhere: their density per unit log-scale is constant. Integrating a constant density over an interval in log-energy gives a linear accumulation. When a coupling's response is the count of resolvable defects, this yields $g^{-2}(t) = g^{-2}(t_0) + \kappa\rho(t - t_0)$, which is precisely the one-loop form with $b = \kappa\rho/2$. The linear-in-log running is therefore a consequence of scale-invariant defect density, not of a one-loop diagram. The sign of asymptotic freedom also emerges: positive density means resolving more structure dilutes the response.

\begin{definition}
The number of defects a probe at log-scale $t$ can resolve, given a scale-invariant density $\rho$ per unit log-scale and reference log-scale $t_0$, is
\begin{equation}
\text{resolvedCount}(\rho, t_0, t) := \rho \cdot (t - t_0) .
\end{equation}
\lean{Running.resolvedCount}
\end{definition}

\begin{theorem}
At the reference log-scale, the resolved count is zero:
\begin{equation}
\text{resolvedCount}(\rho, t_0, t_0) = 0 .
\end{equation}
\lean{Running.resolvedCount\_self}
\end{theorem}

\begin{proof}
Immediate from the definition.
\end{proof}

\begin{definition}
The inverse-square coupling, as accumulated response with each resolvable defect contributing $\kappa$, is
\begin{equation}
u(u_0, \kappa, \rho, t_0, t) := u_0 + \kappa \cdot \text{resolvedCount}(\rho, t_0, t) = u_0 + \kappa \rho (t - t_0) .
\end{equation}
\lean{Running.invSqCoupling}
\end{definition}

\begin{theorem}
At the reference scale, the inverse-square coupling takes its base value:
\begin{equation}
u(u_0, \kappa, \rho, t_0, t_0) = u_0 .
\end{equation}
\lean{Running.invSqCoupling\_at\_base}
\end{theorem}

\begin{proof}
Immediate from the definition.
\end{proof}

\begin{theorem}
The accumulated response has constant slope in log-scale. The function $t \mapsto u(u_0, \kappa, \rho, t_0, t)$ satisfies
\begin{equation}
\frac{du}{dt}(t) = \kappa \rho
\end{equation}
for all $t \in \mathbb{R}$. Nothing perturbative has been used: the slope is a density times a per-defect response.
\lean{Running.hasDerivAt\_invSqCoupling}
\end{theorem}

\begin{proof}
The derivative of $u_0 + \kappa \rho(t - t_0)$ with respect to $t$ is $\kappa \rho$ by the rules of differentiation.
\end{proof}

\begin{theorem}
The counting law and the one-loop equation are identical. Writing $b = \kappa\rho/2$, the slope $\kappa\rho$ is exactly $2b$ of the one-loop form. So the beta function's shape is not an input: it is what a constant density of resolvable defects produces.
\lean{Running.slope\_eq\_two\_b}
\end{theorem}

\begin{proof}
Direct algebra: $2 \cdot (\kappa\rho/2) = \kappa\rho$.
\end{proof}

\begin{theorem}
The counting law reproduces the one-loop running. For any $u_0, \kappa, \rho, t_0, t \in \mathbb{R}$,
\begin{equation}
u(u_0, \kappa, \rho, t_0, t) = u_0 + 2 \left(\frac{\kappa\rho}{2}\right) (t - t_0) .
\end{equation}
\lean{Running.invSqCoupling\_eq\_one\_loop}
\end{theorem}

\begin{proof}
Algebraic simplification of the definition.
\end{proof}

\begin{theorem}[Asymptotic freedom as positive density]
Suppose $\kappa > 0$. Then the response $g^2 = 1/g^{-2}$ falls to zero at high log-scale if and only if the accumulated count grows, i.e., if and only if the defect density is positive:
\begin{equation}
\left(\forall \varepsilon > 0, \exists T : \forall t > T, u(u_0, \kappa, \rho, t_0, t) > 1/\varepsilon\right) \iff \rho > 0 .
\end{equation}
No gauge group or diagram appears; what appears is a count. Resolving more structure dilutes the response: this is the physical content of asymptotic freedom in this framework.
\lean{Running.asympt\_free\_iff\_pos\_density}
\end{theorem}

\begin{proof}
We prove both directions.
($\Rightarrow$) Suppose the left-hand side holds but assume $\rho \leq 0$. For any $\varepsilon > 0$, there exists $T$ such that for all $t > T$, $u(u_0, \kappa, \rho, t_0, t) > 1/\varepsilon$. Taking $\varepsilon = 1/(|u_0| + 1)$, for sufficiently large $t$ we have $u(u_0, \kappa, \rho, t_0, t) > |u_0| + 1$. But if $\rho \leq 0$, then the accumulated count is non-positive for $t \geq t_0$, so $u(u_0, \kappa, \rho, t_0, t) \leq u_0 \leq |u_0|$, a contradiction.
($\Leftarrow$) Suppose $\rho > 0$ and $\varepsilon > 0$. Set $T := t_0 + (1/\varepsilon - u_0)/(\kappa\rho) + 1$. Then for $t > T$,
\begin{equation}
u(u_0, \kappa, \rho, t_0, t) = u_0 + \kappa\rho(t - t_0) > u_0 + \kappa\rho\left(\frac{1/\varepsilon - u_0}{\kappa\rho} + 1\right) = u_0 + (1/\varepsilon - u_0) + \kappa\rho > 1/\varepsilon .
\end{equation}
\end{proof}

\subsection{QCD}

The strong interaction reveals dimensional transmutation: a coupling that is a pure number, running under the scale-invariant framework, secretes an invariant scale $\Lambda$ that appears nowhere in the axioms. Scale invariance is not violated by fiat; it is broken by the solutions of a scale-covariant equation. Proved here is the one-loop solution $g^{-2}(t) = g^{-2}(0) + 2bt$ with $b > 0$, asymptotic freedom, the dimensional transmutation constant $\Lambda = \exp(t - g^{-2}/(2b))$, and the consequence that all hadron mass ratios are parameter-free. The simplest geometric tower of defect masses fails against lattice glueball data, falsifying a prior hypothesis; the parameter-free ratio prediction survives.

\begin{theorem}[Exact solution of one-loop running]
Suppose $g : \mathbb{R} \to \mathbb{R}$ is nonzero and satisfies the one-loop equation
\begin{equation}
\frac{dg}{dt}(t) = -b \cdot (g(t))^3
\end{equation}
for some constant $b \in \mathbb{R}$ and all $t \in \mathbb{R}$. Then
\begin{equation}
(g(t))^{-2} = (g(0))^{-2} + 2bt .
\end{equation}
\lean{QCD.running\_inv\_sq}
\end{theorem}

\begin{proof}
Define $u(r) := (g(r))^{-2} - 2br$. We compute
\begin{equation}
\dot{u}(r) = -\frac{2(g(r))^{-3} \cdot (-b(g(r))^3)}{1} - 2b = 2b - 2b = 0 .
\end{equation}
Thus $u$ is constant, so $u(t) = u(0)$, yielding $(g(t))^{-2} = (g(0))^{-2} + 2bt$.
\end{proof}

\begin{theorem}[Asymptotic freedom]
Suppose $b > 0$, $g : \mathbb{R} \to \mathbb{R}$ is nonzero and satisfies $dg/dt = -b g^3$. Then for any $\varepsilon > 0$, there exists $T \in \mathbb{R}$ such that for all $t > T$, $(g(t))^2 < \varepsilon$. The strong sector becomes free precisely where the local scale is small, as demanded by A3.
\lean{QCD.asymptotic\_freedom}
\end{theorem}

\begin{proof}
From the one-loop solution, $(g(t))^{-2} = (g(0))^{-2} + 2bt$. Choose $T$ large enough that $2bT > 1/\varepsilon - (g(0))^{-2}$. Then for $t > T$, $(g(t))^{-2} > 1/\varepsilon$, so $(g(t))^2 < \varepsilon$.
\end{proof}

\begin{definition}
The scale secreted by the running coupling is defined as
\begin{equation}
\Lambda(b, g, t) := \exp\left(t - \frac{(g(t))^{-2}}{2b}\right) .
\end{equation}
\lean{QCD.Lambda}
\end{definition}

\begin{theorem}[Dimensional transmutation]
The quantity $\Lambda(b, g, t)$ is constant along the flow: a theory whose only input is a dimensionless number has produced a scale, and by A5 no choice of fiducial can shift it. This is how a scale-invariant axiom system comes to describe a world with definite hadron masses. The invariance is not broken in the axioms; it is broken by the solution.
\begin{equation}
\Lambda(b, g, t) = \Lambda(b, g, 0)
\end{equation}
for all $t \in \mathbb{R}$, provided $b \neq 0$.
\lean{QCD.lambda\_invariant}
\end{theorem}

\begin{proof}
From the one-loop solution, $(g(t))^{-2} - (g(0))^{-2} = 2bt$, so
\begin{equation}
\frac{(g(t))^{-2}}{2b} - t = \frac{(g(0))^{-2}}{2b} - t + \frac{(g(0))^{-2} - (g(t))^{-2}}{2b} + t = \frac{(g(0))^{-2}}{2b} - t .
\end{equation}
Thus $t - (g(t))^{-2}/(2b) = -($ constant $)$, so $\Lambda(b, g, t) = \Lambda(b, g, 0)$.
\end{proof}

\begin{theorem}
The strong scale is genuinely positive: $\Lambda(b, g, t) > 0$ for all $b \in \mathbb{R}$, $g : \mathbb{R} \to \mathbb{R}$, and $t \in \mathbb{R}$.
\lean{QCD.lambda\_pos}
\end{theorem}

\begin{proof}
The exponential function is always positive.
\end{proof}

\begin{definition}
A hadron mass is defined as
\begin{equation}
M_{\text{hadron}}(c, \Lambda) := c \cdot \Lambda
\end{equation}
where $c$ is a pure (dimensionless) number and $\Lambda$ is the transmuted scale.
\lean{QCD.hadronMass}
\end{definition}

\begin{theorem}[Parameter-free mass ratios]
The ratio of any two hadron masses is a ratio of pure numbers, independent of the transmuted scale $\Lambda$. Hence it is independent of every input to the theory:
\begin{equation}
\frac{M_{\text{hadron}}(c_1, \Lambda)}{M_{\text{hadron}}(c_2, \Lambda)} = \frac{c_1}{c_2}
\end{equation}
provided $c_2 \neq 0$ and $\Lambda \neq 0$. This is the sharpest prediction the framework makes about the strong sector. Glueballs are the right place to test it because their masses involve no quark mass, so $\Lambda$ is genuinely the only scale.
\lean{QCD.hadron\_mass\_ratio}
\end{theorem}

\begin{proof}
Direct division: $(c_1 \Lambda) / (c_2 \Lambda) = c_1/c_2$.
\end{proof}

\begin{theorem}
Rescaling $\Lambda$ moves every mass together: the spectrum has one overall normalisation and no other freedom. If $k \neq 0$,
\begin{equation}
M_{\text{hadron}}(c, k\Lambda) = k \cdot M_{\text{hadron}}(c, \Lambda) .
\end{equation}
\lean{QCD.spectrum\_rigid}
\end{theorem}

\begin{proof}
$M_{\text{hadron}}(c, k\Lambda) = c(k\Lambda) = k(c\Lambda) = k \cdot M_{\text{hadron}}(c, \Lambda)$.
\end{proof}

\begin{definition}
The ratio of the first two lattice glueball masses is
\begin{equation}
r_1 := \frac{2390}{1710} \approx 1.398 .
\end{equation}
The ratio of the next two is
\begin{equation}
r_2 := \frac{2560}{2390} \approx 1.071 .
\end{equation}
\lean{QCD.ratio$_1$} and \lean{QCD.ratio$_2$}
\end{definition}

\begin{theorem}[The geometric tower fails]
The two consecutive ratios differ by more than $0.3$, so a single geometric tower (which would predict equal ratios) is excluded by lattice data:
\begin{equation}
|r_1 - r_2| > 0.3 .
\end{equation}
This refutes the hypothesis that the defect masses follow the posited law of Defect.mass, not the defect picture itself. A spectrum with strain energy quadratic in the winding, for instance, is unaffected.
\lean{QCD.geometric\_tower\_fails}
\end{theorem}

\begin{proof}
Numerical verification: $|1.398 - 1.071| \approx 0.327 > 0.3$.
\end{proof}

\begin{theorem}
The numerical values of the ratios are:
\begin{enumerate}
\item $1.39 < r_1 < 1.40$.
\item $1.07 < r_2 < 1.08$.
\end{enumerate}
\lean{QCD.ratio$_1$\_value} and \lean{QCD.ratio$_2$\_value}
\end{theorem}

\begin{proof}
Direct arithmetic from the definitions.
\end{proof}

\begin{definition}
With the transmuted scale $\Lambda \approx 332$ MeV, the lightest glueball (at $\approx 1710$ MeV) sits at
\begin{equation}
\frac{m_0}{\Lambda} \approx \frac{1710}{332} \approx 5.15 .
\end{equation}
This is a pure number the framework must eventually compute rather than fit.
\lean{QCD.lightestGlueballInLambda}
\end{definition}

\begin{theorem}
The numerical value is bracketed:
\begin{equation}
5.1 < \frac{m_0}{\Lambda} < 5.2 .
\end{equation}
\lean{QCD.lightestGlueballInLambda\_value}
\end{theorem}

\begin{proof}
Direct arithmetic from the definition.
\end{proof}

\subsection{Spectrum}

The framework predicts the structure of threshold placement. A5 says no log-scale is preferred, so thresholds cannot bunch anywhere: they are placed at a constant rate $\rho$ per unit log-scale. A constant-rate point process has exponentially distributed gaps, whose squared coefficient of variation is exactly one. This prediction discriminates sharply against the geometric tower (CV = 0) and is tested against the eleven consecutive gaps in $\ln m$ across the twelve distinct Standard Model mass scales. The observed $CV^2 = 0.875$ (so $CV = 0.935$) sits at the 63rd percentile of the prediction's sampling range and is incompatible with the geometric tower.

\begin{definition}
The threshold of a structure of mass $m > 0$ is the log-energy at which it becomes resolvable:
\begin{equation}
\mu(m) := \ln m .
\end{equation}
Conversely, a structure of log-threshold $\mu$ has mass $e^\mu$.
\lean{Spectrum.thresholdOfMass} and \lean{Spectrum.massOfThreshold}
\end{definition}

\begin{theorem}
The threshold and mass are inverse maps: $\mu(e^\mu) = \mu$ for all $\mu \in \mathbb{R}$.
\lean{Spectrum.mass\_threshold\_inverse}
\end{theorem}

\begin{proof}
$\mu(e^\mu) = \ln(e^\mu) = \mu$.
\end{proof}

\begin{definition}
For a list $\mathbf{x} = (x_1, \ldots, x_n)$ of real numbers, the \textbf{mean} is
\begin{equation}
\bar{x} := \frac{1}{n}\sum_{i=1}^n x_i ,
\end{equation}
and the \textbf{population variance} is
\begin{equation}
\sigma^2 := \frac{1}{n}\sum_{i=1}^n x_i^2 - \bar{x}^2 .
\end{equation}
\lean{Spectrum.mean} and \lean{Spectrum.variance}
\end{definition}

\begin{theorem}[Geometric tower has zero variance]
If every gap equals a constant $c$, the mean is $c$ and the mean square is $c^2$, so the variance vanishes. A geometric spectrum therefore predicts $CV = \sigma/\bar{x} = 0$, with no freedom:
\begin{equation}
\text{variance}([c, c, \ldots, c]) = 0 .
\end{equation}
\lean{Spectrum.equal\_gaps\_variance\_zero}
\end{theorem}

\begin{proof}
If all $n$ entries are $c$, then $\bar{x} = c$ and $\overline{x^2} = c^2$, so $\sigma^2 = c^2 - c^2 = 0$.
\end{proof}

\begin{definition}
The gaps (in $\ln m$) between consecutive Standard Model mass scales are (taking massless states and unknown neutrino masses into account):
\begin{equation}
\text{gaps} = [1.46, 0.759, 2.985, 0.128, 2.486, 0.336, 0.855, 2.956, 0.126, 0.317, 0.321] .
\end{equation}
The sum is $\sum \text{gaps} = 12.729$ and the sum of squares is $\sum (\text{gaps})^2 = 27.615749$.
\lean{Spectrum.gaps}, \lean{Spectrum.gapSum}, and \lean{Spectrum.gapSqSum}
\end{definition}

\begin{definition}
The mean gap is $\langle \text{gap} \rangle := 12.729/11 \approx 1.1572$.
The population variance of the gaps is $\sigma_{\text{gap}}^2 := 27.615749/11 - (12.729/11)^2$.
The squared coefficient of variation is $CV^2 := \sigma_{\text{gap}}^2 / \langle \text{gap} \rangle^2$.
\lean{Spectrum.meanGap}, \lean{Spectrum.gapVariance}, and \lean{Spectrum.cvSq}
\end{definition}

\begin{definition}
The predictions are:
\begin{enumerate}
\item \textbf{Scale-invariant}: $CV^2 = 1$ (exponential gap distribution).
\item \textbf{Geometric tower}: $CV^2 = 0$ (equal gaps).
\end{enumerate}
\lean{Spectrum.predictedCvSq} and \lean{Spectrum.geometricCvSq}
\end{definition}

\begin{theorem}[Observed coefficient of variation]
The observed value is bracketed:
\begin{equation}
0.87 < CV^2 < 0.88 ,
\end{equation}
i.e., $CV^2 = 0.875$, so $CV \approx 0.935$. This is close to but below the predicted value of $1$.
\lean{Spectrum.observed\_cvSq\_value}
\end{theorem}

\begin{proof}
Direct arithmetic from the gap data.
\end{proof}

\begin{theorem}
The mean gap is bracketed:
\begin{equation}
1.15 < \langle \text{gap} \rangle < 1.16 .
\end{equation}
The inverse is the threshold density $\rho \approx 0.86$ per e-fold.
\lean{Spectrum.meanGap\_value}
\end{theorem}

\begin{proof}
Direct arithmetic from the gap sum.
\end{proof}

\begin{theorem}[The Standard Model spectrum is not geometric]
The gap variance is positive, whereas a geometric tower has variance exactly zero. The observed spectrum is therefore incompatible with the geometric tower:
\begin{equation}
CV^2 \neq 0 .
\end{equation}
This is an independent and sharper refutation than the glueball ratio test, using the whole observed spectrum rather than three lattice numbers.
\lean{Spectrum.observed\_not\_geometric}
\end{theorem}

\begin{proof}
We have $CV^2 \approx 0.875 > 0$, contradicting the geometric prediction of $CV^2 = 0$.
\end{proof}

\begin{theorem}
The observed value is compatible with the scale-invariant prediction. The central 90\% of the sampling distribution for eleven gaps runs from approximately $0.33$ to $1.64$ in $CV^2$:
\begin{equation}
0.33 < CV^2 < 1.64 .
\end{equation}
The observed value $0.875$ sits at the 63rd percentile.
\lean{Spectrum.observed\_within\_scale\_invariant\_range}
\end{theorem}

\begin{proof}
Direct inequalities from the observed value.
\end{proof}

\begin{definition}
Given the mean gap, the threshold density is recovered as
\begin{equation}
\rho(\langle \text{gap} \rangle) := \frac{1}{\langle \text{gap} \rangle} .
\end{equation}
\lean{Spectrum.densityFromGaps}
\end{definition}

\begin{theorem}
With the observed mean gap, the threshold density satisfies
\begin{equation}
0.86 < \rho < 0.87 .
\end{equation}
\lean{Spectrum.densityFromGaps\_value}
\end{theorem}

\begin{proof}
Direct arithmetic from the mean gap.
\end{proof}

\begin{definition}
From the dark-energy equation of state parameter $w$, the dissipation exponent is
\begin{equation}
q(w) := \frac{3w + 5}{2} .
\end{equation}
With $w = -1$, we get $q = 1$.
\lean{Spectrum.qFromW}
\end{definition}

\begin{definition}
From the asymptotic Hubble rate $\lambda = H_\infty$, the dark-energy density is
\begin{equation}
\rho_\Lambda(\lambda, G) := \frac{3\lambda^2}{8\pi G} .
\end{equation}
\lean{Spectrum.rhoLambdaFromLambda}
\end{definition}

\begin{definition}
When a threshold appears, the newly resolvable content is characterized by the jump in the running coefficient $\Delta b$. The increment to the particle count is
\begin{equation}
\Delta N(\Delta b, \kappa) := \frac{\Delta b}{\kappa} .
\end{equation}
\lean{Spectrum.deltaNFromDeltaB}
\end{definition}

\begin{theorem}
The identity $\Delta N \cdot \kappa = \Delta b$ holds (under suitable non-degeneracy conditions).
\lean{Spectrum.deltaN\_inverse}
\end{theorem}

\begin{proof}
Direct algebra: $({\Delta b}/{\kappa}) \cdot \kappa = \Delta b$.
\end{proof}

\begin{theorem}
The conversions are mutually consistent: a spectrum with mean gap $g$ has density $1/g$, and a density $\rho$ has mean gap $1/\rho$.
\begin{equation}
\frac{1}{\rho(g)} = g
\end{equation}
provided $g \neq 0$.
\lean{Spectrum.density\_gap\_inverse}
\end{theorem}

\begin{proof}
$1/(\rho(g)) = 1/(1/g) = g$.
\end{proof}

\subsection{Entropy}

Nothing in the microscopic dynamics distinguishes a direction of time: the substrate is reversible. What breaks the symmetry is that observation happens at a scale, which imposes a finite resolution. A finite resolution is a non-injective map from micro-configurations to what is observable. Under reversible microscopic dynamics followed by coarse-graining, the number of micro-configurations compatible with observation never decreases (the H-theorem), and entropy production is exact bookkeeping, not a statistical assumption. If entropy strictly increases once, then no entropy-non-decreasing law can undo the evolution, making the macroscopic equation irreversible even though the microscopic one is not—precisely the status of the Boltzmann equation.

\begin{definition}
Given a finite set $X$ of micro-configurations and a resolution map $\pi : X \to Y$, the \textbf{saturation} or \textbf{macrostate} generated by a subset $S \subseteq X$ is
\begin{equation}
\text{sat}(\pi, S) := \{x \in X : \pi(x) \in \pi(S)\} .
\end{equation}
It is the set of all micro-configurations that the resolution cannot distinguish from something in $S$.
\lean{Entropy.sat}
\end{definition}

\begin{theorem}
A configuration is in the saturation if and only if its image under $\pi$ lies in the image of $S$:
\begin{equation}
x \in \text{sat}(\pi, S) \iff \pi(x) \in \pi(S) .
\end{equation}
\lean{Entropy.mem\_sat}
\end{theorem}

\begin{proof}
Immediate from the definition.
\end{proof}

\begin{theorem}
Coarse-graining never loses a configuration: $S \subseteq \text{sat}(\pi, S)$.
\lean{Entropy.subset\_sat}
\end{theorem}

\begin{proof}
If $x \in S$, then $\pi(x) \in \pi(S)$, so $x \in \text{sat}(\pi, S)$.
\end{proof}

\begin{theorem}
Therefore, $|S| \leq |\text{sat}(\pi, S)|$.
\lean{Entropy.card\_le\_card\_sat}
\end{theorem}

\begin{proof}
The cardinality of a subset is at most that of the whole set.
\end{proof}

\begin{theorem}[Microscopic reversibility]
If $T : X \to X$ is a bijection (permutation), then any finite subset $S \subseteq X$ has $|T(S)| = |S|$. Microscopic evolution preserves cardinality exactly.
\lean{Entropy.card\_image\_perm}
\end{theorem}

\begin{proof}
The image of a finite set under an injective map has the same cardinality.
\end{proof}

\begin{definition}
One step of observed evolution is: evolve microscopically by $T : X \to X$, then coarse-grain by the resolution $\pi : X \to Y$:
\begin{equation}
F(\pi, T, S) := \text{sat}(\pi, T(S)) .
\end{equation}
\lean{Entropy.evolve}
\end{definition}

\begin{theorem}[H-theorem]
Observed entropy—the log of the number of micro-configurations compatible with the observation—never decreases. Under reversible microscopic dynamics followed by coarse-graining,
\begin{equation}
|S| \leq |F(\pi, T, S)|
\end{equation}
for any subset $S \subseteq X$ and bijection $T : X \to X$. Nothing statistical is invoked.
\lean{Entropy.entropy\_nondecreasing}
\end{theorem}

\begin{proof}
We have $|S| = |T(S)|$ by microscopic reversibility, and $|T(S)| \leq |\text{sat}(\pi, T(S))|$ by coarse-graining. Thus $|S| \leq |F(\pi, T, S)|$.
\end{proof}

\begin{theorem}[The arrow of time]
Let $F, G : \mathcal{P}(X) \to \mathcal{P}(X)$ be any two evolution laws both obeying the H-theorem, i.e., $|S| \leq |G(S)|$ for all $S$. If $F$ strictly produces entropy somewhere, i.e., there exists $S_0$ with $|S_0| < |F(S_0)|$, then there is no inverse law with $G(F(S)) = S$ for all $S$:
\begin{equation}
\neg \exists G : \forall S, |S| \leq |G(S)| \text{ and } G(F(S)) = S \text{ and } \exists S_0, |S_0| < |F(S_0)| .
\end{equation}
So the macroscopic equation cannot be run backwards by a law of its own kind. Time-reversal invariance is lost at the observed level while holding exactly at the microscopic level—which is the precise sense in which the Boltzmann equation is irreversible although the mechanics beneath it is not.
\lean{Entropy.no\_reverse\_law}
\end{theorem}

\begin{proof}
Suppose such a $G$ existed with $G(F(S_0)) = S_0$ and $|F(S_0)| > |S_0|$. By the H-theorem for $G$, $|F(S_0)| \leq |G(F(S_0))| = |S_0|$, contradicting $|F(S_0)| > |S_0|$.
\end{proof}

\begin{theorem}
The same conclusion applies to the concrete observed dynamics with resolution $\pi$ and microscopic evolution $T$.
\lean{Entropy.evolve\_no\_reverse\_law}
\end{theorem}

\begin{proof}
Apply the previous theorem to $F = F(\pi, T)$.
\end{proof}

\begin{theorem}
In equilibrium every observable is stationary. Any quantity that changes under the dynamics therefore witnesses departure from equilibrium—Sakharov's third condition, stated contrapositively.
\lean{Entropy.charge\_change\_implies\_nonequilibrium}
\end{theorem}

\begin{proof}
If $F(S) = S$ (equilibrium) but $Q(F(S)) \neq Q(S)$, then $Q(S) \neq Q(S)$, a contradiction.
\end{proof}

\begin{theorem}
Strict entropy production is itself a departure from equilibrium. With the H-theorem above, an open dissipating universe satisfies Sakharov's third condition at every moment at which it produces entropy.
\lean{Entropy.entropy\_production\_implies\_nonequilibrium}
\end{theorem}

\begin{proof}
If $|S| < |F(S)|$ (entropy production) but $F(S) = S$ (equilibrium), then $|S| < |S|$, a contradiction.
\end{proof}

\subsection{Light}

The question ``what is light?'' has a sharp answer in this framework. The measured quadratic form is $Q_g = s^2 Q_\delta$, where $s$ is a unit. Multiplying by a unit cannot create or destroy a zero. Therefore: the null cone is scale-blind; light does not see the scale field $\sigma$ at all. For every non-null direction the scale is recoverable from the ratio of the measured to bare form—matter sees the scale and is the only thing that does. Light is the scale-blind part of the geometry, weight-zero under scale transformations, which is why it alone can carry the comparison of local units from one event to another. The null condition is therefore invariant under the fiducial shift of A5, and the scale connection is flat (by the conformal part of the framework).

\begin{definition}
The \textbf{bare quadratic form} of the substrate, with signature $\eta : \text{Fin } n \to A$ (entries $\pm 1$), is
\begin{equation}
Q_{\text{bare}}(\eta, v) := \sum_{i=0}^{n-1} \eta(i) \cdot v(i)^2 .
\end{equation}
This is pure bookkeeping and carries no physics.
\lean{Light.bareForm}
\end{definition}

\begin{definition}
The \textbf{measured quadratic form}, read in the local unit $s \in A^\times$, is
\begin{equation}
Q_{\text{meas}}(s, \eta, v) := s^2 \cdot Q_{\text{bare}}(\eta, v) .
\end{equation}
This is A4 applied to a tangent vector.
\lean{Light.physForm}
\end{definition}

\begin{definition}
A direction $v$ is \textbf{null} if the measured form vanishes on it: $Q_{\text{meas}}(s, \eta, v) = 0$.
\lean{Light.IsNull}
\end{definition}

\begin{theorem}
Multiplying by a unit neither creates nor destroys a zero. For any unit $s \in A^\times$ and element $x \in A$,
\begin{equation}
s^2 \cdot x = 0 \iff x = 0 .
\end{equation}
\lean{Light.unit\_sq\_mul\_eq\_zero\_iff}
\end{theorem}

\begin{proof}
If $s^2 \cdot x = 0$, then $x = (s^2)^{-1} \cdot (s^2 \cdot x) = 0$. If $x = 0$, then $s^2 \cdot x = 0$ trivially.
\end{proof}

\begin{theorem}[The null cone is scale-blind]
The set of null directions is determined by the bare bookkeeping alone. No choice of scale field $\sigma$—that is, no gravitational field—can alter which directions are null:
\begin{equation}
Q_{\text{meas}}(s, \eta, v) = 0 \iff Q_{\text{bare}}(\eta, v) = 0 .
\end{equation}
\lean{Light.isNull\_iff\_bare}
\end{theorem}

\begin{proof}
By the preceding theorem, $s^2 \cdot Q_{\text{bare}}(\eta, v) = 0$ if and only if $Q_{\text{bare}}(\eta, v) = 0$.
\end{proof}

\begin{theorem}[Causal structure is scale-independent]
Two observers using any two scale fields $s$ and $s'$ agree exactly about which directions are null:
\begin{equation}
Q_{\text{meas}}(s, \eta, v) = 0 \iff Q_{\text{meas}}(s', \eta, v) = 0 .
\end{equation}
\lean{Light.isNull\_scale\_invariant}
\end{theorem}

\begin{proof}
Both are equivalent to $Q_{\text{bare}}(\eta, v) = 0$ by the previous theorem.
\end{proof}

\begin{theorem}[The scale is measured by massive probes]
If the bare form on $v$ is a regular element (i.e., $\forall x : A, x \cdot Q_{\text{bare}}(\eta, v) = 0 \to x = 0$), and two observers with scales $s$ and $s'$ assign the same measured form to $v$,
\begin{equation}
Q_{\text{meas}}(s, \eta, v) = Q_{\text{meas}}(s', \eta, v) ,
\end{equation}
then $s^2 = (s')^2$. Combined with the scale-blindness of the null cone, this is the complete epistemology: light fixes the conformal class and nothing more; everything about the scale must be read off with matter.
\lean{Light.scale\_determined\_of\_nonnull}
\end{theorem}

\begin{proof}
We have $(s^2 - (s')^2) \cdot Q_{\text{bare}}(\eta, v) = Q_{\text{meas}}(s, \eta, v) - Q_{\text{meas}}(s', \eta, v) = 0$. Since $Q_{\text{bare}}(\eta, v)$ is regular, $s^2 - (s')^2 = 0$, so $s^2 = (s')^2$.
\end{proof}

\begin{theorem}
On a non-null direction, the measured and bare forms differ unless the scale is trivial. If $Q_{\text{bare}}(\eta, v)$ is regular and $s^2 \neq 1$, then
\begin{equation}
Q_{\text{meas}}(s, \eta, v) \neq Q_{\text{bare}}(\eta, v) .
\end{equation}
\lean{Light.physForm\_ne\_bareForm}
\end{theorem}

\begin{proof}
If they were equal, then $(s^2 - 1) \cdot Q_{\text{bare}}(\eta, v) = 0$. Since $Q_{\text{bare}}(\eta, v)$ is regular, $s^2 - 1 = 0$, contradicting $s^2 \neq 1$.
\end{proof}

\begin{theorem}[Light has weight zero]
The null condition is invariant under any change of scale, in particular under the fiducial shift of A5. A global rescaling $s \to us$ leaves null directions null:
\begin{equation}
Q_{\text{meas}}(s, \eta, v) = 0 \iff Q_{\text{meas}}(us, \eta, v) = 0 .
\end{equation}
Light is exactly the weight-zero part of the geometry, which is the formal content of masslessness.
\lean{Light.null\_weight\_zero}
\end{theorem}

\begin{proof}
Both are equivalent to $Q_{\text{bare}}(\eta, v) = 0$ by scale-blindness.
\end{proof}

\subsection{Observation}

An internal tension in the development is resolved here. Light is scale-blind, but observers are made of matter, which carries the scale. A null vector has no length, but it has energy—assigned by an observer. Redshift is not a property of light; it is the ratio of two observers' scales. This resolves the apparent conflict: what is measured in cosmology is not the global drift (which A5 declares unobservable) but the scale difference between emission and absorption.

\begin{definition}
The \textbf{pairing} of a probe direction $v$ with an observer's direction $w$, in bare bookkeeping, is
\begin{equation}
\langle w, v \rangle_{\text{bare}}(\eta, w, v) := \sum_{i=0}^{n-1} \eta(i) \cdot w(i) \cdot v(i) .
\end{equation}
Unlike the bare form, this is not a property of $v$ alone; it requires the observer.
\lean{Observation.pairing}
\end{definition}

\begin{definition}
The \textbf{energy an observer assigns} to $v$, read in that observer's local unit $u \in A^\times$, is
\begin{equation}
E_{\text{obs}}(u, \eta, w, v) := u^2 \cdot \langle w, v \rangle_{\text{bare}}(\eta, w, v) .
\end{equation}
Energy is a two-place quantity; nothing about $v$ by itself has an energy.
\lean{Observation.obsEnergy}
\end{definition}

\begin{theorem}[A null direction carries energy]
Being null is a scale-blind statement about $v$ alone. Having energy is a statement about $v$ and an observer, which is not scale-blind. The two are consistent because they are not about the same thing. Explicitly, there exist a signature $\eta$, a null direction $v$, and an observer direction $w$ such that $Q_{\text{bare}}(\eta, v) = 0$ but $\langle w, v \rangle_{\text{bare}} \neq 0$.
\lean{Observation.null\_has\_energy}
\end{theorem}

\begin{proof}
Take the signature $\eta = (-1, 1)$ on $\mathbb{F}^2_2 = \mathbb{Z}^2$. Let $v = (1, 1)$ and $w = (1, 0)$. Then $Q_{\text{bare}} = (-1) \cdot 1 + 1 \cdot 1 = 0$, and $\langle w, v \rangle_{\text{bare}} = (-1) \cdot 1 \cdot 1 + 1 \cdot 0 \cdot 1 = -1 \neq 0$.
\end{proof}

\begin{theorem}[Redshift is a ratio of observer scales]
Two observers with local units $u_1$ and $u_2$ assign to the same null direction energies whose ratio is $(u_1/u_2)^2$. The light contributed nothing; the whole ratio is the difference between the two local units:
\begin{equation}
E_{\text{obs}}(u_1, \eta, w, v) \cdot u_2^2 = E_{\text{obs}}(u_2, \eta, w, v) \cdot u_1^2 .
\end{equation}
This is why an unobservable global drift nonetheless produces an observable redshift: what is measured is the scale difference between emission and absorption, which A5 declares physical.
\lean{Observation.redshift\_is\_scale\_ratio}
\end{theorem}

\begin{proof}
We compute $u_1^2 \langle w, v \rangle_{\text{bare}} \cdot u_2^2 = u_2^2 \langle w, v \rangle_{\text{bare}} \cdot u_1^2$, which is algebraic.
\end{proof}

\begin{theorem}[No redshift without scale difference]
If two observers assign energies to the same null direction whose cross-products differ, their units differ:
\begin{equation}
E_{\text{obs}}(u_1, \eta, w, v) \cdot u_2^2 \neq E_{\text{obs}}(u_2, \eta, w, v) \cdot u_1^2 \implies u_1 \neq u_2 .
\end{equation}
Equal units give equal energies by congruence; that is $x = x$, not a theorem about redshift.
\lean{Observation.scale\_differs\_of\_energy\_differs}
\end{theorem}

\begin{proof}
Contrapositive of the redshift theorem: if $u_1 = u_2$, the cross-products are equal.
\end{proof}

\begin{theorem}
Rescaling both observers' units by the same factor leaves the energy ratio untouched. If $c \in A^\times$,
\begin{equation}
E_{\text{obs}}(cu, \eta, w, v) = c^2 \cdot E_{\text{obs}}(u, \eta, w, v) .
\end{equation}
Rescaling both changes nothing observable—this is A5.
\lean{Observation.obsEnergy\_rescale}
\end{theorem}

\begin{proof}
$E_{\text{obs}}(cu, \eta, w, v) = (cu)^2 \langle w, v \rangle_{\text{bare}} = c^2 u^2 \langle w, v \rangle_{\text{bare}} = c^2 E_{\text{obs}}(u, \eta, w, v)$.
\end{proof}

\begin{theorem}[A5 is respected]
A global rescaling applied to both observers leaves the observed ratio untouched. For any $c, u_1, u_2 \in A^\times$,
\begin{equation}
E_{\text{obs}}(cu_1, \eta, w, v) \cdot (cu_2)^2 = E_{\text{obs}}(cu_2, \eta, w, v) \cdot (cu_1)^2 .
\end{equation}
\lean{Observation.redshift\_fiducial\_invariant}
\end{theorem}

\begin{proof}
Apply the redshift theorem to $cu_1$ and $cu_2$.
\end{proof}

\begin{theorem}
The causal structure remains scale-blind. For any two scale fields $s$ and $s'$ in a ScaleAlgebra, the null condition for a direction $v$ is independent of which scale field is used:
\begin{equation}
Q_{\text{meas}}(s, \eta, v) = 0 \iff Q_{\text{meas}}(s', \eta, v) = 0 .
\end{equation}
The theorems above about energy do not weaken this; energy simply is not the null condition.
\lean{Observation.causal\_structure\_still\_blind}
\end{theorem}

\begin{proof}
This restates the scale-blindness of Light.isNull\_iff\_bare.
\end{proof}

\subsection{Particles}

Taking the scale–energy duality A3 ($\varepsilon \cdot s = 1$) seriously for excitations yields two consequences. An excitation of log-scale $\mu$ (i.e., log-energy $-\mu$) is a configuration whose own energy scale is $e^{-\mu}$. It can be produced only where that energy is available, so it does not exist as a degree of freedom below its threshold—the particle content is a function of scale, and effective field theory reports its own structure. A decaying excitation carries its own clock, set by its rest energy $m$. The laboratory uses a different unit, set by total energy $E = \sqrt{p^2 + m^2}$. A lifetime is measured in both units, giving $\gamma = E/m$; time dilation is the scale ratio of A3, read twice. The reported lifetime is at least the proper one and grows with momentum.

\begin{definition}
The total energy of an excitation of rest scale $m$ carrying momentum $p$ is
\begin{equation}
E(m, p) := \sqrt{p^2 + m^2} .
\end{equation}
All quantities are dimensionless, measured against the same fiducial.
\lean{Particles.energy}
\end{definition}

\begin{definition}
The scale ratio between the laboratory's energy scale and the excitation's own is
\begin{equation}
\gamma(m, p) := \frac{E(m, p)}{m} .
\end{equation}
This is the only content of the Lorentz factor in this framework.
\lean{Particles.scaleRatio}
\end{definition}

\begin{theorem}
For $m > 0$, the energy is positive: $E(m, p) > 0$ for all $p \in \mathbb{R}$.
\lean{Particles.energy\_pos}
\end{theorem}

\begin{proof}
$E(m, p) = \sqrt{p^2 + m^2} \geq \sqrt{m^2} = m > 0$.
\end{proof}

\begin{theorem}
For $m > 0$, the energy is at least the rest scale: $E(m, p) \geq m$ for all $p \in \mathbb{R}$.
\lean{Particles.energy\_ge}
\end{theorem}

\begin{proof}
$\sqrt{p^2 + m^2} \geq \sqrt{m^2} = m$.
\end{proof}

\begin{theorem}
The scale ratio is at least one. The laboratory never assigns an excitation a smaller scale than its own rest scale:
\begin{equation}
\gamma(m, p) \geq 1
\end{equation}
for all $m > 0$ and $p \in \mathbb{R}$.
\lean{Particles.one\_le\_scaleRatio}
\end{theorem}

\begin{proof}
$\gamma(m, p) = E(m, p)/m \geq m/m = 1$.
\end{proof}

\begin{theorem}
The scale ratio is monotone in momentum. If $p^2 \leq q^2$, then $\gamma(m, p) \leq \gamma(m, q)$ for any $m > 0$.
\lean{Particles.scaleRatio\_mono}
\end{theorem}

\begin{proof}
Both $E(m, p) \leq E(m, q)$, so $\gamma(m, p) \leq \gamma(m, q)$.
\end{proof}

\begin{definition}
The lifetime reported by a laboratory is the excitation's proper lifetime $\tau_0$, re-expressed in the laboratory's unit and thus multiplied by the scale ratio:
\begin{equation}
\tau(u_0, m, p) := \gamma(m, p) \cdot \tau_0 .
\end{equation}
\lean{Particles.lifetime}
\end{definition}

\begin{theorem}
The reported lifetime is never shorter than the proper one:
\begin{equation}
\tau_0 \leq \tau(u_0, m, p)
\end{equation}
for all $m > 0$ and $\tau_0 \geq 0$.
\lean{Particles.lifetime\_ge}
\end{theorem}

\begin{proof}
$\tau = \gamma \cdot \tau_0 \geq 1 \cdot \tau_0 = \tau_0$ since $\gamma \geq 1$.
\end{proof}

\begin{theorem}
Lifetime grows with momentum purely because the scale ratio grows:
\begin{equation}
p^2 \leq q^2 \implies \tau(u_0, m, p) \leq \tau(u_0, m, q)
\end{equation}
for all $m > 0$ and $\tau_0 \geq 0$. No dynamical assumption about the decay is used.
\lean{Particles.lifetime\_mono}
\end{theorem}

\begin{proof}
Both $\tau(u_0, m, p) = \gamma(m, p) \tau_0 \leq \gamma(m, q) \tau_0 = \tau(u_0, m, q)$.
\end{proof}

\begin{theorem}
At rest the two scales coincide and the reported lifetime is the proper one:
\begin{equation}
\tau(u_0, m, 0) = \tau_0 .
\end{equation}
The normalisation is fixed with nothing left over.
\lean{Particles.lifetime\_at\_rest}
\end{theorem}

\begin{proof}
$\gamma(m, 0) = \sqrt{0 + m^2}/m = 1$, so $\tau = 1 \cdot \tau_0 = \tau_0$.
\end{proof}

\begin{definition}
An excitation of log-scale $\mu$ is \textbf{available} at log-energy $t$ if $\mu \leq t$.
\lean{Particles.Available}
\end{definition}

\begin{theorem}
Availability is monotone: if an excitation is available at $t$ and $t \leq t'$, then it is available at $t'$.
\lean{Particles.available\_mono}
\end{theorem}

\begin{proof}
$\mu \leq t \leq t'$ implies $\mu \leq t'$.
\end{proof}

\begin{theorem}
Below its threshold an excitation is not a degree of freedom at all. If $t < \mu$, then the excitation is not available at $t$.
\lean{Particles.not\_available\_of\_lt}
\end{theorem}

\begin{proof}
Immediate from the definition of availability.
\end{proof}

\subsection{DarkMatter}

The framework forbids the usual question ``what particle is it?'' and replaces it with a question about scales. A detector at energy scale $\varepsilon_D$ has spatial resolution $1/\varepsilon_D$; what it can register is an interaction that moves its own state by at least its threshold $\theta$. An excitation whose energy scale $\varepsilon$ is far below $\varepsilon_D$ deposits an arbitrarily small fraction of the detector's scale—the interaction happens, the observation fails. Gravity couples to total energy $N \cdot \varepsilon$, which is blind to how that energy is partitioned. A large multiplicity of very-low-scale excitations is gravitationally loud and non-gravitationally silent—dark matter's entire observational profile.

\begin{definition}
The fraction of a detector's own energy scale that a single excitation of energy scale $\varepsilon$ can move is
\begin{equation}
R_{\text{det}}(\varepsilon, \varepsilon_D) := \frac{\varepsilon}{\varepsilon_D} .
\end{equation}
This is what a detector actually measures.
\lean{Dark.detResp}
\end{definition}

\begin{definition}
The gravitational response, which depends on total energy and is blind to how it is partitioned, is
\begin{equation}
R_{\text{grav}}(N, \varepsilon) := N \cdot \varepsilon ,
\end{equation}
where $N$ is the multiplicity.
\lean{Dark.gravResp}
\end{definition}

\begin{theorem}
The detector response is positive for positive $\varepsilon$ and $\varepsilon_D$:
\begin{equation}
\varepsilon > 0, \varepsilon_D > 0 \implies R_{\text{det}}(\varepsilon, \varepsilon_D) > 0 .
\end{equation}
\lean{Dark.detResp\_pos}
\end{theorem}

\begin{proof}
Positive numbers divide to a positive number.
\end{proof}

\begin{theorem}[Interaction without observation]
For any detector scale $\varepsilon_D > 0$ and detection threshold $\theta > 0$, there exists an energy scale $\varepsilon_0$ such that every excitation of scale $\varepsilon < \varepsilon_0$ interacts ($R_{\text{det}} > 0$) yet falls below threshold ($R_{\text{det}} < \theta$). The failure is in the measurement, not in the coupling; the probe and apparatus live at incommensurable scales.
\lean{Dark.detection\_failure}
\end{theorem}

\begin{proof}
Choose $\varepsilon_0 := \varepsilon_D \cdot \theta$. For $0 < \varepsilon < \varepsilon_0$, we have $R_{\text{det}}(\varepsilon, \varepsilon_D) = \varepsilon/\varepsilon_D > 0$ and $R_{\text{det}}(\varepsilon, \varepsilon_D) = \varepsilon/\varepsilon_D < \varepsilon_0/\varepsilon_D = \theta$.
\end{proof}

\begin{theorem}[Dark matter exists]
Fix any gravitational signal $M > 0$, detector scale $\varepsilon_D > 0$, and detection threshold $\theta > 0$. Then there exist an energy scale $\varepsilon > 0$ and a multiplicity $N > 0$ such that $R_{\text{grav}}(N, \varepsilon) = M$ while $R_{\text{det}}(\varepsilon, \varepsilon_D) < \theta$. Gravitational visibility and non-gravitational invisibility are not in tension; they are two different functionals of the same configuration, separated by a scale hierarchy.
\lean{Dark.dark\_matter\_exists}
\end{theorem}

\begin{proof}
Let $\varepsilon_0 = \varepsilon_D \theta$ from the preceding theorem. Choose $\varepsilon := \min(\varepsilon_0/2, M) > 0$ (which is positive by the hypotheses). Set $N := M/\varepsilon > 0$. Then $R_{\text{grav}}(N, \varepsilon) = (M/\varepsilon) \varepsilon = M$, and since $\varepsilon < \varepsilon_0$, we have $R_{\text{det}}(\varepsilon, \varepsilon_D) < \theta$.
\end{proof}

\begin{theorem}
The detectable response of a fixed gravitational mass tends to zero as the constituent scale is lowered. For any $\varepsilon_D > 0$ and $\theta > 0$, there exists $\varepsilon_0 > 0$ such that for all $\varepsilon \in (0, \varepsilon_0)$,
\begin{equation}
R_{\text{det}}(\varepsilon, \varepsilon_D) < \theta .
\end{equation}
\lean{Dark.response\_vanishes}
\end{theorem}

\begin{proof}
This is the limit form of the detection failure theorem.
\end{proof}

\subsection{DarkEnergy}

The universe is not closed. Axiom A6 states that the total dimensionless energy scale $\varepsilon$ dissipates. By the scale–energy duality A3 ($\varepsilon \cdot s = 1$), a decaying energy scale is a growing length scale. Cosmic expansion is not motion through space; it is the running down of the energy scale, read through the duality. We prove: constant-rate dissipation $d\varepsilon/dt = -\lambda \varepsilon$ has unique solution $\varepsilon(t) = \varepsilon_0 e^{-\lambda t}$; the scale factor is exact de Sitter $a(t) = a_0 e^{\lambda t}$; the Hubble rate is constant $H = \lambda$; and the expansion accelerates $\ddot{a} = \lambda^2 a > 0$. Dark energy is not a substance with negative pressure; it is the dissipation rate of an open universe, and the observed $\Lambda$ is $\lambda^2$ up to dimensionless factors.

\begin{theorem}[Auxiliary: derivative of exponential]
For any $\lambda, t \in \mathbb{R}$,
\begin{equation}
\frac{d}{du} e^{\lambda u}\bigg|_{u=t} = \lambda e^{\lambda t} .
\end{equation}
\lean{Cosmology.hasDerivAt\_exp\_mul}
\end{theorem}

\begin{proof}
Chain rule: $\frac{d}{du} e^{\lambda u} = e^{\lambda u} \cdot \lambda$.
\end{proof}

\begin{theorem}[Dissipation has unique history]
If the dimensionless energy scale obeys $d\varepsilon/dt = -\lambda \varepsilon$, then $\varepsilon(t) = \varepsilon(0) e^{-\lambda t}$. Nothing is assumed about the form of $\varepsilon$; the exponential is forced.
\begin{equation}
\varepsilon(t) = \varepsilon(0) \cdot e^{-\lambda t} .
\end{equation}
\lean{Cosmology.dissipation\_solution}
\end{theorem}

\begin{proof}
Define $u(r) := \varepsilon(r) e^{\lambda r}$. Then
\begin{equation}
\dot{u}(r) = (-\lambda \varepsilon(r)) e^{\lambda r} + \varepsilon(r) \cdot \lambda e^{\lambda r} = 0 .
\end{equation}
Thus $u$ is constant, so $\varepsilon(t) e^{\lambda t} = \varepsilon(0)$, yielding $\varepsilon(t) = \varepsilon(0) e^{-\lambda t}$.
\end{proof}

\begin{theorem}
The energy scale of an open universe is strictly decreasing. If $\lambda > 0$, $\varepsilon(0) > 0$, and $t > 0$, then
\begin{equation}
\varepsilon(t) < \varepsilon(0) .
\end{equation}
\lean{Cosmology.energy\_strictly\_decreasing}
\end{theorem}

\begin{proof}
From the solution, $\varepsilon(t) = \varepsilon(0) e^{-\lambda t}$ with $e^{-\lambda t} < 1$ for $t > 0$ and $\lambda > 0$.
\end{proof}

\begin{definition}
The cosmological scale factor, defined by duality from the energy scale, is
\begin{equation}
a(a_0, \lambda, t) := a_0 \cdot e^{\lambda t} .
\end{equation}
\lean{Cosmology.scaleFactor}
\end{definition}

\begin{theorem}[Dissipation is expansion]
The reciprocal of a decaying energy scale is an exponentially growing length scale: exact de Sitter, with no cosmological constant inserted anywhere. If $\varepsilon(0) \neq 0$,
\begin{equation}
\frac{1}{\varepsilon(t)} = a\left(\frac{1}{\varepsilon(0)}, \lambda, t\right) = \frac{1}{\varepsilon(0)} \cdot e^{\lambda t} .
\end{equation}
\lean{Cosmology.scaleFactor\_eq}
\end{theorem}

\begin{proof}
From $\varepsilon(t) = \varepsilon(0) e^{-\lambda t}$, we have $1/\varepsilon(t) = (1/\varepsilon(0)) e^{\lambda t}$.
\end{proof}

\begin{theorem}[Hubble rate is constant]
The Hubble rate $H = \dot{a}/a$ is constant and equal to the dissipation rate:
\begin{equation}
\frac{d}{dt} a(a_0, \lambda, t) = \lambda \cdot a(a_0, \lambda, t) .
\end{equation}
Thus $H = \lambda$.
\lean{Cosmology.hasDerivAt\_scaleFactor}
\end{theorem}

\begin{proof}
$\frac{d}{dt}(a_0 e^{\lambda t}) = a_0 \lambda e^{\lambda t} = \lambda \cdot a(a_0, \lambda, t)$.
\end{proof}

\begin{theorem}[Accelerated expansion]
The expansion accelerates at a rate proportional to the square of the dissipation rate:
\begin{equation}
\frac{d^2}{dt^2} a(a_0, \lambda, t) = \lambda^2 \cdot a(a_0, \lambda, t) > 0
\end{equation}
for any nonzero dissipation rate and positive scale. Acceleration is not evidence of a repulsive substance; it is the second derivative of a reciprocal.
\lean{Cosmology.hasDerivAt\_hubble}
\end{theorem}

\begin{proof}
$\frac{d}{dt}(\lambda a_0 e^{\lambda t}) = \lambda^2 a_0 e^{\lambda t} = \lambda^2 a(a_0, \lambda, t)$. Since $a > 0$ and $\lambda \neq 0$, this is positive.
\end{proof}

\begin{theorem}
The second derivative of the scale factor is strictly positive for any $a_0 > 0$ and $\lambda \neq 0$:
\begin{equation}
\ddot{a}(a_0, \lambda, t) = \lambda^2 \cdot a(a_0, \lambda, t) > 0 .
\end{equation}
\lean{Cosmology.acceleration\_pos}
\end{theorem}

\begin{proof}
We have $a(a_0, \lambda, t) = a_0 e^{\lambda t} > 0$ since $a_0 > 0$. Thus $\lambda^2 a > 0$.
\end{proof}

\begin{theorem}[Collected statement: open universe with dissipation]
An open universe dissipating at a constant fractional rate $\lambda > 0$ with initial energy scale $\varepsilon_0 > 0$ expands, and does so at an accelerating rate, with constant Hubble parameter $H = \lambda$. More precisely, for any $\varepsilon : \mathbb{R} \to \mathbb{R}$ satisfying $d\varepsilon/dt = -\lambda \varepsilon$ with $\varepsilon(0) = \varepsilon_0$:
\begin{enumerate}
\item $\displaystyle \frac{1}{\varepsilon(t)} = \frac{1}{\varepsilon_0} \cdot e^{\lambda t}$ for all $t \in \mathbb{R}$.
\item $\displaystyle \frac{d}{dt} a\left(\frac{1}{\varepsilon_0}, \lambda, t\right) = \lambda \cdot a\left(\frac{1}{\varepsilon_0}, \lambda, t\right)$ for all $t \in \mathbb{R}$.
\item $\displaystyle 0 < \lambda^2 \cdot a\left(\frac{1}{\varepsilon_0}, \lambda, t\right)$ for all $t \in \mathbb{R}$.
\end{enumerate}
\lean{Cosmology.open\_universe\_accelerates}
\end{theorem}

\begin{proof}
(1) Follows from the dissipation solution with $\varepsilon(0) = \varepsilon_0$.
(2) The derivative of $a(1/\varepsilon_0, \lambda, t) = (1/\varepsilon_0) e^{\lambda t}$ is $\lambda \cdot (1/\varepsilon_0) e^{\lambda t} = \lambda \cdot a(1/\varepsilon_0, \lambda, t)$.
(3) Since $a > 0$ and $\lambda \neq 0$ (implied by $\lambda > 0$), the product is positive.
\end{proof}

